%% ICAIF '26 submission — ACM sigconf, double-blind (anonymous).
%% Compile on Overleaf (acmart is built in) or with: pdflatex -> bibtex -> pdflatex x2.
%% Hard limit: 8 pages including figures AND references. No appendix/supplementary.
%% Remove the `anonymous` option for the camera-ready to reveal the author block below.
\documentclass[sigconf,anonymous,review]{acmart}

%% --- ACM boilerplate: no copyright block for review; set real values at camera-ready. ---
\settopmatter{printacmref=false}
\renewcommand\footnotetextcopyrightpermission[1]{}
\pagestyle{plain}
\setcopyright{none}
\acmConference[ICAIF '26]{ACM International Conference on AI in Finance}{November 14--17, 2026}{Milan, Italy}

\usepackage{booktabs}
\usepackage{amsmath}

\begin{document}

\title{Skill or Structure? Auditing Learned Skill and\\ Explanation Faithfulness in Residual Reinforcement Learning\\ for Portfolio Allocation}

%% Author block is suppressed by the `anonymous` class option for double-blind review.
\author{Meghanadh}
\affiliation{%
  \institution{Independent Researcher}
  \country{}
}
\email{meghanadh27@gmail.com}

\renewcommand{\shortauthors}{Anonymous}

\begin{abstract}
Reinforcement learning (RL) allocators are often credited with reducing
out-of-sample left-tail risk, yet long-only allocation already inherits
downside protection for free from the simplex constraint, log-reward geometry,
and the classical rebalancing premium. We ask how much of an RL allocator's
apparent skill is genuinely learned and how much is inherited structure, and
whether post-hoc explanations of the trained policy identify its true decision
mechanism. We train PPO as a \emph{residual} on a structural-null base policy
(equal-weight, volatility-scaled, or risk-parity) with a
\emph{structure-baselined} reward that credits only return earned above the
base, and we calibrate the resulting skill measure on a synthetic market with a
toggleable predictive signal: skill rises monotonically with forecastability and
vanishes when prediction is impossible. On real ETFs (2005--2024, 10\,bps costs,
walk-forward), both a scalar de-risking gate and a more expressive per-asset tilt
add \emph{no} skill above structure once measured net of a phase-randomization
placebo null; the expressive agent loses more. We then probe the trained policy
with per-feature causal ablation and compare it to gradient saliency and
KernelSHAP. Attribution is broadly faithful to the causal mechanism, but which
method is most faithful is agent-dependent, and SHAP recurrently over-credits a
salient, high-variance feature. The synthetic method-agreement structure
transfers to the real agent.
\end{abstract}

\begin{CCSXML}
<ccs2012>
 <concept><concept_id>10010147.10010257.10010258.10010261.10010269</concept_id>
 <concept_desc>Computing methodologies~Reinforcement learning</concept_desc>
 <concept_significance>500</concept_significance></concept>
 <concept><concept_id>10003752.10010070.10010071</concept_id>
 <concept_desc>Theory of computation~Sequential decision making</concept_desc>
 <concept_significance>300</concept_significance></concept>
 <concept><concept_id>10010147.10010178.10010179</concept_id>
 <concept_desc>Computing methodologies~Machine learning</concept_desc>
 <concept_significance>300</concept_significance></concept>
</ccs2012>
\end{CCSXML}
\ccsdesc[500]{Computing methodologies~Reinforcement learning}
\ccsdesc[300]{Theory of computation~Sequential decision making}

\keywords{reinforcement learning, portfolio allocation, explainability,
attribution, causal intervention, tail risk, residual policy learning}

\maketitle

\section{Introduction}
Model-free reinforcement learning has been reported to beat mean--variance
optimization and reduce out-of-sample left-tail risk in portfolio allocation
\cite{lavko2023,zhang2020dlportfolio}. Two long-standing critiques temper this
claim. First, evaluations frequently omit transaction costs and rarely clear a
strong $1/N$ benchmark \cite{demiguel2009,kruthof2025}. Second, the ``skill''
being measured may not be skill at all: a long-only allocator inherits downside
protection mechanically from the probability-simplex constraint, from
maximizing log wealth, and from the rebalancing premium of Stochastic Portfolio
Theory \cite{fernholz2002}. A volatility-managed or risk-parity rule delivers
much of the reported tail benefit with no forecasting \cite{moreira2017,maillard2010}.

We take these critiques as a measurement problem and ask two questions.
\textbf{(RQ-skill)} How much of an RL allocator's apparent tail-risk benefit is
learned skill above the free structure, rather than the structure itself?
\textbf{(RQ-faithful)} When we explain the trained policy with standard
attribution tools, do those explanations identify the mechanism the policy
actually uses?

Our approach isolates skill by construction. We fix a \emph{structural-null base
policy} that already harvests the free structure, train the agent as a
\emph{residual} on that base, and reward it only for return earned \emph{above}
the base on the same market path. A zero residual reproduces the base exactly, so
the zero-skill floor is free and the gradient never sees the structural growth.
We calibrate this skill measure on a synthetic market whose predictive signal can
be switched off, giving a ground-truth boundary where genuine skill is impossible.
We then apply the calibrated measure to real ETFs, and finally probe the trained
policy's decision mechanism with causal interventions, adjudicating gradient
saliency and KernelSHAP against what those interventions reveal.

\noindent\textbf{Contributions.}
(1) A structure-baselined residual-RL skill measure, validated against synthetic
ground truth where skill rises with forecastability and vanishes without it.
(2) A real-data verdict: neither a scalar de-risking gate nor an expressive
per-asset tilt adds skill above structure net of a placebo null, across three
bases and many seeds; the expressive agent loses \emph{more}.
(3) A causal-vs-attribution faithfulness probe across three agents, showing
attribution is broadly faithful but its reliability is agent-dependent, with a
recurring SHAP salience bias, and that the synthetic method-agreement transfers
to the real agent.

We frame success as a rigorous, honest verdict, not as beating the market.

\section{Related Work}
RL allocators span value- and policy-based methods and are increasingly paired
with interpretability claims \cite{lavko2023,cong2021alphaportfolio,delarica2025}.
Recent cost- and benchmark-aware replications question whether the RL advantage
survives realistic frictions and a $1/N$ baseline \cite{kruthof2025,demiguel2009}.
Our residual formulation borrows residual policy learning from robotics, where a
learned policy corrects a fixed controller \cite{silver2018,johannink2019}; we
repurpose it as a measurement instrument, pairing it with a structure-baselined
reward and a synthetic ground-truth test. On explanation, gradient saliency
\cite{simonyan2014} and KernelSHAP \cite{lundberg2017} are the standard tools,
but saliency maps for deep RL have been shown to be exploratory rather than
explanatory \cite{atrey2020}. We contribute a controlled, causally grounded test
of what these tools recover for RL allocators specifically.

\section{Method}

\subsection{Structure-baselined residual RL}
Let $b_t \in \Delta^{n-1}$ be the weights of a fixed base policy on the
probability simplex, computed from a trailing window (equal-weight,
volatility-scaled, or equal-risk-contribution risk parity \cite{maillard2010}).
The agent acts as a residual on $b_t$ in one of two modes.

\emph{De-risking gate.} A scalar action $g_t \in [0,1]$ shifts weight toward a
safe asset $s$: $w_t = (1-g_t)\,b_t + g_t\,e_s$, where $e_s$ is the unit vector on
$s$. $g_t=0$ reproduces the base; larger $g_t$ times timed de-risking.

\emph{Expressive tilt.} A per-asset action $a_t \in \mathbb{R}^n$ produces
$w_t = \Pi_{\Delta}\big(b_t + \tau \tanh(a_t)\big)$, where $\Pi_\Delta$ is
Euclidean projection onto the simplex \cite{duchi2008} and $\tau$ bounds the tilt.
A zero action again reproduces the base, but the agent can express any long-only
deviation.

The reward is \emph{structure-baselined}: the agent's net log return minus the
base's net log return on the same path, each charged its own turnover cost,
\begin{equation}
r_t = \log\!\big(1 + w_t^\top \rho_t - c\,\|w_t - w_{t-1}\|_1\big)
     - \log\!\big(1 + b_t^\top \rho_t - c\,\|b_t - b_{t-1}\|_1\big),
\end{equation}
where $\rho_t$ are asset returns and $c$ the per-unit cost. The cumulative reward
is the agent's realized \emph{skill above structure}; it is zero for the base
itself, so the gradient sees only what the agent adds.

We train PPO \cite{schulman2017} with an MLP policy via Stable-Baselines3
\cite{raffin2021}. The observation is a trailing return window plus realized
volatility, momentum, the base weights, and a de-risking signal.

\subsection{Synthetic ground truth}
To calibrate the skill measure we generate a regime-switching market with a
\emph{toggleable} leading signal of strength $\sigma \in [0,1]$. When $\sigma$ is
high, timed de-risking is possible; when $\sigma=0$ the market is unforecastable
and genuine skill is impossible by construction. A valid skill measure must rise
with $\sigma$ and vanish at $\sigma=0$.

\subsection{Real-data null}
On real data there is no ground truth for ``no skill.'' We therefore report skill
\emph{net of a phase-randomization placebo null}: surrogate return series that
preserve each asset's spectrum but destroy volatility clustering, retrained
identically. The placebo mean is the luck/overfitting floor; a positive floor is
expected and is exactly why netting is necessary. We report skill net of this
null with a small-sample confidence interval and the fraction of placebo draws
that match or beat the real agent (placebo exceedance).

\subsection{Causal-vs-attribution faithfulness probe}
Given a trained policy we replay its deterministic decisions and measure two
tracks over the observation features. The \emph{causal} track ablates each
feature (freeze to its mean; permute across time) and records the mean absolute
change in the policy's scalar decision, a direct measure of what the policy
responds to. The \emph{attribution} track computes gradient saliency
\cite{simonyan2014} in gradient$\times$standard-deviation units and KernelSHAP
\cite{lundberg2017}. Putting saliency in gradient$\times$std units and expressing
all methods as normalized per-feature shares removes a scale confound: the signal
feature has far larger raw scale than the return features, and raw
scale-invariant saliency would otherwise be incomparable to the scale-sensitive
causal and SHAP tracks. We report per-feature Spearman correlation between the
causal profile and each attribution profile. On synthetic data the leading signal
is the known driver; on real data no such oracle exists, so only
\emph{method agreement} is adjudicable, not faithfulness against truth.

\section{Experiments}
The real universe is 11 liquid ETFs spanning equity sectors, Treasuries, gold,
and international equity (2005--2024, daily, 10\,bps costs), evaluated with an
anchored expanding-window walk-forward and per-fold retraining. Synthetic runs
use 150k training steps and five seeds. All figures below are held-out.

\subsection{The skill measure is valid (synthetic)}
Table~\ref{tab:rq2} shows mean skill rising monotonically with signal strength
for the gate, and vanishing toward the unforecastable null. The expressive tilt
cannot reach a costless do-nothing floor (with no signal it over-churns and loses
about $-6.4\times10^{-5}$), so we judge it net of that churn null: skill on minus
skill off is $+2.0\times10^{-4}$, roughly four times the floor. The measure
detects learned skill only when learned skill is possible, for both agents.

\begin{table}[t]
  \caption{Synthetic validation: mean structure-baselined skill (held-out,
  5 seeds). Skill rises with forecastability and vanishes at the null.}
  \label{tab:rq2}
  \small
  \begin{tabular}{lc}
    \toprule
    Signal strength $\sigma$ (gate) & Mean skill \\
    \midrule
    0.00 (unforecastable null) & $5.9\times10^{-5}$ \\
    0.50                       & $1.7\times10^{-4}$ \\
    0.95 (strongly forecastable) & $2.9\times10^{-4}$ \\
    \midrule
    \multicolumn{2}{l}{\emph{Expressive tilt}, net-of-null (signal 0.95):} \\
    \quad skill off / skill on & $-6.4\times10^{-5}$ / $+1.4\times10^{-4}$ \\
    \quad skill net of churn floor & $+2.0\times10^{-4}$ ($4.1\times$) \\
    \bottomrule
  \end{tabular}
\end{table}

\subsection{No skill above structure (real ETFs)}
Table~\ref{tab:rq1} reports skill net of the placebo null across the three bases,
for both agents. Every skill-net confidence interval lies entirely below zero,
and every structureless surrogate beats the real agent (placebo exceedance
$=1.00$) for all bases. The gate learns to mostly do nothing (mean gate
$\approx 0.04$), tracking its base on all tail and risk metrics while both clearly
beat naive $1/N$; the tail benefit over $1/N$ is therefore inherited structure,
not learned skill. The expressive tilt loses \emph{more} than the gate
($1.5$--$1.8\times$), worst for the weakest base: more expressiveness means more
churn, and where no timeable structure exists that churn is pure cost drag. A
capable agent sharpens the verdict rather than rescuing it.

\begin{table}[t]
  \caption{Real ETFs: skill net of the placebo null (5 agent seeds, 10 placebo
  draws per base). All intervals lie below zero; placebo exceedance $=1.00$
  throughout. The expressive tilt loses more than the gate.}
  \label{tab:rq1}
  \small
  \begin{tabular}{lcc}
    \toprule
    Structural base & Gate skill net & Tilt skill net \\
    \midrule
    Equal-weight & $-1.76\times10^{-4}$ & $-2.06\times10^{-4}$ \\
    Vol-scaled   & $-1.19\times10^{-4}$ & $-1.75\times10^{-4}$ \\
    Risk-parity  & $-1.01\times10^{-4}$ & $-1.54\times10^{-4}$ \\
    \bottomrule
  \end{tabular}
\end{table}

\subsection{Is the explanation faithful? (RQ-faithful)}
Table~\ref{tab:rq3} summarizes the probe across three agents. On the synthetic
gate, where the signal is the known driver, KernelSHAP names it the top feature in
all five seeds and gradient saliency in four of five (the miss is the least
signal-concentrated agent); per-feature agreement with the causal profile is
positive for both. On the synthetic expressive tilt, saliency is near-perfectly
aligned with the causal profile while SHAP over-credits the salient, high-variance
signal feature, inverting which method is more reliable. A methodological caveat
appears here: with wildly unequal group cardinalities (many return features versus
one signal), group-level ``top driver'' verdicts are cardinality-confounded, so
the per-feature Spearman is the robust measure.

On the real agent there is no ground truth, so we report method agreement only.
The activity diagnostic overturns a natural prior: although the mean gate is
$\approx 0.05$, its standard deviation is $\approx 0.18$, so the agent actively
swings the gate rather than sitting still. It has a real, if unprofitable,
decision mechanism (none of five seeds degenerate). Given that live mechanism,
per-feature agreement with causal ablation is strong: saliency $+0.996$ and SHAP
$+0.79$, preserving the saliency-over-SHAP ordering seen synthetically. All four
methods attribute about $0.96$ of the normalized share to the recent-returns
block, with the crisis signal a minor contributor; unlike the synthetic tilt this
is not a cardinality artifact, since the signal is only about three times an
average return feature here. The real allocator has a coherent, faithfully
readable decision rule; the rule simply is not skillful.

\begin{table}[t]
  \caption{Faithfulness probe: per-feature Spearman between the causal profile and
  each attribution method (5 seeds). Real data has no ground truth, so its numbers
  are method agreement, not faithfulness.}
  \label{tab:rq3}
  \small
  \begin{tabular}{lccl}
    \toprule
    Agent & Saliency & SHAP & Note \\
    \midrule
    Synthetic gate & $+0.79$ & $+0.61$ & SHAP names driver 5/5 \\
    Synthetic tilt & $+0.999$ & $+0.88$ & SHAP over-credits signal \\
    Real gate      & $+0.996$ & $+0.79$ & agreement only \\
    \bottomrule
  \end{tabular}
\end{table}

\section{Discussion and Limitations}
The three questions close a loop. The measure detects skill when it exists
(synthetic), finds none on real markets net of luck (both agents), and shows that
the trained policy's mechanism is largely readable by post-hoc attribution, with
saliency tracking the causal per-feature profile more tightly than SHAP across all
three agents. The practical failure mode is SHAP's tendency to over-credit a
salient, high-variance feature; which tool to trust is agent-dependent.

Several limitations bound these claims. The synthetic market is deliberately
simple, chosen so that ``no skill'' is well defined rather than realistic. The
real study uses one liquid universe and daily data; different universes,
frequencies, or richer action spaces may behave differently. Attribution on real
data measures agreement, not faithfulness, because no oracle mechanism exists; high
agreement cannot rule out two methods being wrong together, though causal ablation
is the most direct available reference. The per-feature causal vectors were
aggregated to groups rather than persisted in full, which a camera-ready
extension should save to state per-feature dominance directly. Finally, the
placebo null destroys volatility clustering but preserves linear autocorrelation;
alternative nulls are worth reporting.

\section{Conclusion}
Treating RL portfolio allocation as a measurement problem, we separate learned
skill from inherited structure with a structure-baselined residual and a
synthetic ground-truth calibration, and we test whether standard explanations
identify the trained policy's mechanism. On real ETFs, neither a parsimonious nor
an expressive RL allocator earns skill above a transparent structural base once
costs and a luck null are accounted for. The policy's decision rule is
nonetheless coherent and faithfully readable, most reliably by gradient saliency;
KernelSHAP carries a salience bias. Code and configurations will be released
publicly.

%% --- References ---
\bibliographystyle{ACM-Reference-Format}
\bibliography{references}

\end{document}
